\begin{abstract}

Model-heterogeneous federated learning (HFL) enables clients with different architectures to collaborate, but existing corruption-robust HFL methods mainly treat corruption as a nuisance that degrades input quality. We study a controlled failure mode, \textbf{class--corruption entanglement (CLE)}, in which client-specific associations between task classes and corruption operators can turn corruption into a directional class cue. To distinguish this behavior from ordinary corruption difficulty, we construct source-paired operator interventions and introduce \textbf{Directional Shortcut Alignment (DSA)}, which measures whether predictive probability mass moves toward classes associated with the applied operator during training. Under explicit semantic-preservation, pre-registered-binding, and evaluation-isolation assumptions, DSA identifies a binding-direction operator-response contrast. A complementary proxy non-identifiability result shows that improvements in pseudo-environment statistics or generic proxy objectives cannot, without additional assumptions, certify reduced true CLE dependence. In a seed-0 matched HFL-versus-Local factorial, the observed shortcut is predominantly local-first in the controlled setting, motivating a local intervention. We therefore combine a taxonomy-assisted coarse \textbf{Public Environment Witness (PEW)} with \textbf{Balanced Environment Risk (BER)}, which compresses class-conditional pseudo-environment support imbalance without changing the communication protocol. On a held-out binding map over three matched training seeds, PEW+BER reduces pooled DSA relative to ERM, CVaR-DRO, and PEW+GroupDRO while improving pooled task utility. Additional seed-0 experiments replicate DSA suppression across another binding map and a second protocol-matched HFL communication base. Utility improvements are not uniform across bases or clients, and PEW remains taxonomy-assisted.

\end{abstract}
